% ============================================================
% Chapter 3 — Scope Evolution: Proposal to Delivered System
% ============================================================
\chapter{Scope Evolution: From Proposal to Delivered System}
\label{ch:scope}

\roadmap{This chapter reconciles the eighteen deliverables promised in the
project proposal with the system actually delivered. It presents the full
proposed-versus-delivered accounting with a rationale for every departure,
details the five items whose status requires the most careful reporting,
lists the substantial capabilities that were built without having been
proposed, and closes with the revised scope statement. The framing
throughout is \emph{scope evolution, not scope drift}: each departure was a
research-driven decision, and none of the statuses below has been upgraded
beyond what the codebase supports.}

\section{How the scope evolved}
\label{sec:scope-how}

The proposal described an intelligent tutoring system in which content
authoring, tutoring, and a modest sensing deliverable (``WebGazer tracing
eye gaze and facial expressions,'' item~12) would culminate in a user
study with statistical comparison of system-detected and self-reported
states. During development, a sequence of empirical realizations shifted
the project's center of gravity from \emph{content generation} to
\emph{learning-process measurement}. Three realizations were decisive.
First, for any claim about affect or attention to survive examination, the
raw signals would need to be time-aligned, inspectable, and re-codable
after the fact---which forced the construction of the sync-anchor time
base, the session-replay engine, and the retrospective coding surface
(\cref{ch:sensing}), none of which the proposal anticipated. Second,
single-channel facial sensing proved scientifically insufficient on its
own terms: the literature's insistence on multimodality and on validation
against an independent criterion (\cref{sec:bg-webcam}) motivated the
separate AU pipeline, the AOI attention layer, and the emotion self-report
probe. Third, the arrival of capable LLM APIs mid-project made
retrieval-grounded tutoring and strategy interventions far more powerful
than the proposal's ``chatbot linked to LLM'' had envisioned, and the
project invested accordingly. The cost of these investments was paid by
four proposal items (5, 8, 10, 13) that were delivered in reframed or
partial form; the accounting below reports them as such.

\section{Proposed versus delivered: the eighteen items}
\label{sec:scope-table}

\Cref{tab:scope-18} gives the verdict for every proposal item. Verdicts
follow the codebase audit exactly\footnote{\code{docs/capstone-evidence/02_scope_and_deliverables.md}}:
\emph{Delivered} (built as promised), \emph{Exceeded} (built beyond the
proposal), \emph{Partial} (core exists, promised scope incomplete),
\emph{Reframed} (built differently for stated reasons), and
\emph{Unverifiable from code} (a study-conduct item that source code
cannot attest). No status has been upgraded.

\begin{small}
\begin{longtable}{@{}p{0.05\linewidth}p{0.27\linewidth}p{0.17\linewidth}p{0.42\linewidth}@{}}
\caption{Proposal items: proposed scope, delivered status, and rationale.}
\label{tab:scope-18}\\
\toprule
\# & Proposed & Delivered status & Rationale / evidence \\
\midrule
\endfirsthead
\multicolumn{4}{l}{\small\emph{\tablename~\thetable{} (continued)}}\\
\toprule
\# & Proposed & Delivered status & Rationale / evidence \\
\midrule
\endhead
\bottomrule
\endlastfoot
1 & Login pages for instructor and learners via email authentication &
Delivered & Email-or-login-ID authentication with JWT, role guards, and
rate limiting.\footnote{\code{apps/api/src/auth}} \\
2 & RAG pipeline to upload course materials & Exceeded & Hybrid
dense--sparse retrieval with reciprocal rank fusion, optional reranking,
contextual chunk blurbs, and multimodal page-image embedding
(\cref{sec:modes-rag}). \\
3 & LLM course generation with RAG embeddings & Delivered & Course-outline
and per-page generation jobs grounded in retrieved
chunks.\footnote{\code{apps/api/src/course-structure}, \code{apps/api/src/page-content}} \\
4 & Human safeguard edits & Delivered & Draft approve/reject workflow,
page version history with rollback, and a justification-gated publish
override.\footnote{\code{ContentDraft}, \code{PageContentVersion}, \code{apps/api/src/publish-gate}} \\
5 & Learning-path generation with knowledge components using knowledge
graph theory & \textbf{Partial / Reframed} & Genuine KC graph with
prerequisite edges, topological ordering, and cycle detection exists for
the \emph{teacher}; no runtime student-facing adaptive router. Two KC
systems remain unwired (\cref{sec:scope-item5}). \\
6 & LLM evaluation of course content (formatting, accuracy, pedagogy) &
Delivered (nuanced) & Rubric scores for formatting, equations, pedagogy,
and rigor, plus deterministic math checks and auto-fix; ``accuracy'' is
rubric-based quality evaluation, not external fact-checking
(\cref{sec:scope-item6}). \\
7 & Curation of question banks and assessment & Exceeded & Question CRUD
and bulk import, assessment assembly, and an AI
generate--review--approve pipeline with Hattie-model
grading.\footnote{\code{apps/api/src/questions}, \code{apps/api/src/question-generation}} \\
8 & LMS monitoring: enrolment, progress, performance & \textbf{Partial /
Reframed} & Enrollment policy, mastery views, and deep per-session
analytics exist; the at-a-glance cohort dashboard is a stub
(\cref{sec:scope-item8}). \\
9 & Lesson content page (modular lesson pages) & Delivered & Module/item
structure with block-based pages and PDF lessons rendered in the student
course view.\footnote{\code{ModuleItem}; \code{apps/web/src/pages/student/StudentCourseViewPage.tsx}} \\
10 & Assessment page with adaptive feedback & \textbf{Reframed} &
Adaptivity delivered as feedback \emph{depth} (four Hattie levels,
per-course configuration), not adaptive \emph{item selection}; BKT columns
exist but BKT is not implemented (\cref{sec:scope-item10}). \\
11 & Chatbot linked to LLM for dialogue-based learning & Exceeded & Two
grounded surfaces: the dialogue workspace over learner-uploaded sources
and the docked/floating course chatbot, sharing one grounding contract
(\cref{ch:modes}). \\
12 & WebGazer tracing eye gaze and facial expressions & Exceeded &
WebGazer gaze at 10\,Hz plus two server-side facial pipelines
(OpenFace~3.0 emotions; Py-Feat action units), time-aligned and
replayable (\cref{ch:sensing}). \\
13 & Learners' progress report & \textbf{Partial / Reframed} &
Session summaries, results pages, and mastery views exist; no
consolidated longitudinal per-student report artifact
(\cref{sec:scope-item13}). \\
14 & User study with $\geq$30 participants & Unverifiable from code &
Study conduct is not attestable from source; the codebase references 54
recorded sessions predating the self-report probe. Participant counts are
reported in \cref{ch:results}. \todo{Confirm final N.} \\
15 & Validation of system-detected vs.\ self-reported states & Partial &
Both sides of the comparison exist in the data (self-report probe;
windowed state estimates); the agreement analysis itself is performed in
\cref{ch:results}, not in-platform. \\
16 & T-test and one-way ANOVA (phase 1 vs.\ phase 2) & Partial &
Dependent variables are exportable; \emph{no phase marker exists in the
database}, so grouping must be defined externally
(\cref{sec:methodology-stats}). \\
17 & Mann--Whitney test if normality fails & Partial & Same data basis
and same grouping gap as item 16; normality screening is an analysis-side
step. \\
18 & Report of effect sizes & Partial & All raw quantities exist; no
effect-size or normalized-gain fields are stored, so effect sizes are
computed in analysis (\cref{ch:results}). \\
\end{longtable}
\end{small}

\section{The five items requiring careful reporting}
\label{sec:scope-five}

\subsection{Item 5: knowledge-component learning paths}
\label{sec:scope-item5}

The knowledge-graph deliverable is substantially, but asymmetrically,
built. On the teacher side there is a complete knowledge-component (KC)
studio: AI-proposed KCs, prerequisite/related edges, Kahn topological
ordering and depth-first cycle detection over the graph, KC-to-page
mappings, versioned snapshots, and a learning-path
panel.\footnote{\code{apps/api/src/kc} (proposed-kcs, kc-graph,
kc-mappings); \code{apps/web/src/components/LearningPathPanel.tsx}}
``Knowledge graph theory'' is therefore genuinely implemented. What does
not exist is the student-facing consequence the proposal implied: a
runtime router that selects each learner's next content from mastery over
the prerequisite graph. The structural reason is an architectural seam:
the platform contains \emph{two} KC systems---the curated graph
(\dbtable{ProposedKC}/\dbtable{KcEdge}) and the mastery tracker
(\dbtable{KnowledgeComponent}/\dbtable{UserMastery}, an exponential
moving average with $\alpha=0.3$)---that are not connected. Until mastery
is computed over the graph's nodes, an adaptive path cannot be driven by
it. The report accordingly classifies item 5 as partial/reframed and
identifies the unification as future work.

\subsection{Item 6: LLM evaluation of content}
\label{sec:scope-item6}

Content evaluation is delivered: evaluation runs score each page on
formatting, equations, pedagogy, and rigor, attach located issues with
suggested fixes, combine an LLM rubric pass with a deterministic
mathematical-notation normalizer, and can apply fixes back into the page
under version control.\footnote{\code{apps/api/src/evaluation};
\dbtable{PageEvalResult}} The nuance the report preserves is semantic:
the proposal's word ``accuracy'' is implemented as \emph{rubric-scored
rigor}, i.e., internal quality review, not verification against an
external ground-truth corpus. The report therefore uses
``content-quality evaluation'' and never claims factual verification.

\subsection{Item 8: LMS monitoring}
\label{sec:scope-item8}

Monitoring was reframed from breadth to depth. The delivered system has
enrollment management with per-course policy flags, per-topic mastery
views for students and teachers, per-session behavioral summaries, and
the review hub whose replay, conversation, intervention, and biometrics
tabs constitute an unusually deep per-learner microscope
(\cref{ch:sensing}). What it lacks is the conventional wide-angle lens:
the teacher landing dashboard renders placeholder statistics with no data
fetch,\footnote{\code{apps/web/src/pages/teacher/TeacherDashboard.tsx};
four stat cards hardcoded to ``--''.} and there is no cohort-level
``who is falling behind'' roster. For a research deployment with a small
cohort the depth-first trade was correct; for the proposal's LMS wording
it is a partial delivery, and the report says so.

\subsection{Item 10: adaptive feedback}
\label{sec:scope-item10}

The assessment surface delivers rich feedback---automatic MCQ grading,
LLM grading of short answers, and structured feedback at up to four
levels (task, process, self-regulation, self) following
\citet{hattie2007}, with teachers configuring which levels are enabled
per course. What was \emph{not} built is adaptivity in the
item-selection sense: no computerized-adaptive testing, no difficulty
branching. The mastery model's schema carries Bayesian Knowledge Tracing
parameter columns \citep{corbett1995}, but they are unpopulated; mastery
is an EMA. The closest delivered analogue to adaptive selection is the
revision-worksheet generator, which selects questions in the learner's
zone-of-proximal-development band. The honest formulation, used
throughout this report: adaptivity was delivered in feedback
\emph{depth}, not item \emph{selection}.

\subsection{Item 13: learners' progress report}
\label{sec:scope-item13}

Progress information exists in several partial surfaces---per-session
summaries with mastery deltas, a student results page, dashboard mastery
bars, and exhaustive researcher-facing exports---but no consolidated,
longitudinal ``progress report'' document per learner. The capability
was effectively reframed into the research exports (per-session and
cohort CSV bundles, \cref{sec:sensing-export}); a student- or
parent-facing report generator remains future work.

\section{Deliverables built beyond the proposal}
\label{sec:scope-new}

\Cref{tab:scope-new} lists capabilities present in the delivered system
that the proposal did not anticipate. They are not incidental: the first
five constitute the platform's research instrumentation and are the basis
of contributions 2--6 in \cref{sec:intro-contributions}.

\begin{table}[htbp]
\centering
\small
\caption{Substantial capabilities delivered beyond the proposal's scope.}
\label{tab:scope-new}
\begin{tabular}{@{}p{0.3\linewidth}p{0.62\linewidth}@{}}
\toprule
Capability & Summary \\
\midrule
Session replay engine & Time-aligned replay of screen state, webcam
video, gaze, affect estimates, and events; three-stage loading for long
sessions (\cref{sec:sensing-replay}). \\
AOI/SEEV attention layer & Interface-region gaze bucketing, epoch
segmentation, and the observed-vs-expected \allocscore{}
(\cref{sec:sensing-aoi}). \\
Retrospective coding & Researcher-private code books and time-anchored
annotations over replays (\cref{sec:sensing-coding}). \\
EF text-mining pipeline & Nine construct detectors over learner
utterances with calibration metadata (\cref{sec:methodology-ef}). \\
Emotion self-report probe & Mandatory five-option affect report every 15
minutes; the study's self-report criterion
(\cref{sec:methodology-selfreport}). \\
Affective-mapping engine & Teacher-configurable, versioned rules mapping
emotion probabilities to four learning-centered states
(\cref{sec:sensing-affect}). \\
Sync-anchor time alignment & Wall/monotonic/server clock anchor unifying
all modalities on one time base (\cref{sec:sensing-logging}). \\
Dialogue mode (Mode A) & Source-grounded dialogue workspace with studio
outputs (briefings, flashcards, comparisons, FAQs)
(\cref{sec:modes-a}). \\
Four interventions & Practice testing, interrogative elaboration,
stepwise learning, distributed practice (\cref{ch:interventions}). \\
Teacher prompt configuration & Per-course system prompts for
interventions, feedback levels, and text-mining constructs
(\cref{sec:interventions-teacher}). \\
CSV/ZIP research exporters & Wide-format replay CSV, full-session media
bundle, and a seven-file cohort export (\cref{sec:sensing-export}). \\
Supporting infrastructure & Publish gate, curriculum-coverage analysis,
knowledge versioning, pre-generated exercises, LLM cost accounting,
bulk user provisioning, VLM slide description. \\
\bottomrule
\end{tabular}
\end{table}

\section{Revised scope statement}
\label{sec:scope-statement}

The delivered system is best described as a \emph{multimodal
learning-analytics platform with an intelligent-tutoring front end}. The
authoring side meets, and in places exceeds, the proposal: multimodal RAG
with hybrid retrieval, LLM course and question generation under
human-in-the-loop review, a knowledge-component graph with prerequisite
reasoning, and automated content-quality evaluation. The learning side
adds four theory-grounded, learner-initiated strategy interventions and
two grounded conversational surfaces. Crucially, the sensing side is
decoupled from the pedagogy: interventions are never triggered by
affective signals, so the affect measurements form a clean observational
layer rather than a confound inside the learning loop. On top sits the
instrumentation---synchronized interaction, gaze, pupil-proxy, emotion,
and action-unit streams on a single wall-clock anchor---and the research
toolset (replay, AOI attention scoring, executive-function text mining,
qualitative coding, and exporters) that turns each session into an
analyzable multimodal record.

Several proposal deliverables were reframed rather than dropped, and this
report treats them as such: knowledge-graph ``learning paths'' exist as
teacher-facing prerequisite graphs rather than a runtime adaptive router;
``adaptive feedback'' was delivered as adaptive feedback depth rather
than adaptive item selection; LMS ``monitoring'' and the ``progress
report'' were reframed into fine-grained per-session analytics rather
than a conventional dashboard. The statistical program (validation
against self-report; parametric or non-parametric group comparison;
effect sizes) is supported by the collected data but executed in
analysis rather than in-platform, with two acknowledged gaps---no stored
phase marker and no stored gain fields---handled explicitly in
\cref{ch:methodology}. Framed this way, the platform's evolution is
itself the contribution: a working demonstration that a RAG-based tutor
can be instrumented, end to end and privacy-first, to make learners'
attention, affect, and strategy use measurable.

\medskip
\noindent The next chapter documents the architecture that carries this
scope.
